\section{File Format}\label{sec:problem-file-format}

Problem files are text files, so they can be created with any standard text editor,
such as UNIX \texttt{vi} or the Silicon Graphics \texttt{jot} editor. They can also
be created by other software, such as an interactive trajectory-design program.

Information in a problem file is read into TAOS one line at a time. Each line
contains names and values separated by special characters such as spaces or commas.
These special characters are called \emph{delimiters}. Valid delimiters are
white-space characters (space, tab, and so forth), commas, equal signs,
parentheses, colons, greater-than signs, and less-than signs. Because delimiters are
used to separate names and values, they cannot be embedded within a name or value.

Internally, TAOS requires lowercase names. Names and keywords entered in uppercase
are automatically converted to lowercase, so from a user's standpoint case is not
significant except in titles.

\taossubhead{Names and Values}

Names and values can be placed anywhere on a line; they do not have to be in
particular columns. This free-field style permits indentation that makes the
problem-file structure easier to visualize and helps avoid input errors.

Because input is free-field, the order of names and values is often important. For
example, in the trajectory data block

\begin{lstlisting}[style=taosmanual]
*trajectory 3 missile start on 23
\end{lstlisting}

the keyword \problemblock{*trajectory} must be followed by a trajectory number,
then by a trajectory or vehicle name, the keywords \texttt{start on}, and finally the
starting segment number. In other data blocks, where lists of variables and values
are given, order may be less critical.

Numerical values can be entered as integers without a decimal point, real numbers
with a decimal point, or real numbers in scientific notation (\emph{e}-format). For
example,

\begin{lstlisting}[style=taosmanual]
*initial geodetic
  time=20, vel=2.0e4, wt=345.5, gama=-24.3
\end{lstlisting}

uses \texttt{2.0e4}, which is interpreted as $2.0\times10^4$, or 20,000.
Delimiters cannot be embedded within a value, so

\begin{lstlisting}[style=taosmanual]
wt=10,000       # Incorrect value
\end{lstlisting}

is incorrect. The comma is interpreted as a delimiter, and TAOS reads two values,
10 and 0, instead of 10000.

\taossubhead{Survey, Search, and Optimization Parameters}

One of the powerful features of TAOS is its ability to vary input parameters
systematically for tradeoff studies. Surveys, searches, and optimization loops all
change numerical values in the problem file. Surveys vary parameters through a set
of known values and compute a trajectory for each value. Searches vary one or two
input parameters to achieve specified conditions during the trajectory. For example,
the initial flight-path angle can be varied until the final range equals a desired value.
Optimization loops vary a set of parameters to maximize or minimize a trajectory
quantity subject to constraints. For example, an angle-of-attack history can be varied
to maximize range while maintaining an impact velocity above a required value and
a dynamic pressure above a specified minimum.

All three methods adjust numerical input values. They cannot be used to change
table names or other problem-file keywords.

A numerical field is designated for a survey by assigning it a special keyword. For
example,

\begin{lstlisting}[style=taosmanual]
*initial geodetic
  time=20, vel=2.0e4, wt=surv-1, gama=-24.3
\end{lstlisting}

sets the initial weight to values supplied by survey loop 1 in the corresponding
\problemblock{*survey} block (\cref{sec:block-survey}). Values are designated for
searches with \texttt{srch-n}, where \texttt{n} is the search number
(\cref{sec:block-search}). Optimization values use \texttt{optx-n}, where
\texttt{x} identifies the optimization loop and \texttt{n} identifies the parameter
(\cref{sec:block-optimize}). At this point it is only necessary to know that the
capability exists; the separate survey, search, and optimization sections give the
complete details.

\taossubhead{Titles}

Some data blocks have titles. Titles are special because delimiters are ignored. A
title extends from the first nonblank character to the end of the line. Titles are also
not converted to lowercase; case is preserved as entered.

\taossubhead{Comments}

Comments begin with the special character \texttt{\#} and continue to the end of the
line. They can be placed anywhere in a problem file. Blank lines are ignored and can
be used to separate groups of information for readability. For example,

\begin{lstlisting}[style=taosmanual]
*define alt-km                         # Altitude in km
  alt-km = alt / 3280.84;              # Convert units of altitude
\end{lstlisting}

defines a new variable called \texttt{alt-km} containing vehicle altitude in
kilometres rather than feet. The comments document both the variable and the unit
conversion.
